\documentclass[11pt,a4paper]{article}
\usepackage[utf8]{inputenc}
\usepackage[T1]{fontenc}
\usepackage{amsmath,amsfonts,amssymb}
\usepackage{graphicx}
\usepackage{hyperref}
\usepackage{listings}
\usepackage{color}
\usepackage{booktabs}
\usepackage{multirow}
\usepackage{float}
\usepackage{geometry}
\usepackage{multicol}
\geometry{margin=1in}

% Color definitions
\definecolor{zenblue}{RGB}{41,121,255}
\definecolor{zengreen}{RGB}{52,199,89}
\definecolor{zenorange}{RGB}{255,149,0}
\definecolor{zenpurple}{RGB}{175,82,222}

\hypersetup{
    colorlinks=true,
    linkcolor=zenblue,
    urlcolor=zenblue,
    citecolor=zenblue
}

\title{
    \vspace{-2cm}
    \Huge \textbf{The Zen AI Model Family} \\
    \vspace{0.5cm}
    \Large An Honest Overview: Finetunes and Packaging of Open Models \\
    \vspace{0.3cm}
    \normalsize Technical Overview and Provenance Whitepaper v2.0
}

\author{
    Zach Kelling\thanks{Corresponding author: zach@lux.network} \\
    \textit{Hanzo Industries \quad Lux Industries \quad Zoo Labs Foundation} \\
    \texttt{research@lux.network}
}

\date{September 2025}

\begin{document}

\maketitle

\begin{abstract}
The \textbf{Zen AI Model Family} is a curated set of models for language, vision, and speech that are
\emph{packaged and lightly fine-tuned derivatives of open, permissively licensed upstream models}---
predominantly Alibaba's \textbf{Qwen3} series (Apache-2.0), with a few members based on other open
releases (e.g., DeepSeek and IBM Granite). Zen is not a from-scratch model program: we do not run
large-scale pretraining, and we do not use a proprietary ``Zen MoDE'' architecture. Earlier versions
of this overview claimed homegrown architectures, parameter counts (such as a ``480B'' code model and
``235B'' design models), order-of-magnitude efficiency advantages (``10$\times$ smaller,'' ``up to
98\% less energy''), and per-model benchmark and energy tables. Those claims were not substantiated by
the shipped artifacts and have been removed. This corrected overview states, for each Zen model, the
real upstream base and its license, and defers to the upstream technical reports for pretraining
details and benchmark numbers. The value Zen provides is selection, instruction/format tuning, and
convenient multi-format deployment artifacts (SafeTensors, GGUF, MLX), all under the upstream
licenses with attribution.
\end{abstract}

\tableofcontents
\newpage

\section{Introduction}

The cost of pretraining a competitive foundation model from scratch is out of reach for most teams.
The Zen family takes the pragmatic, now-standard route: start from strong, permissively licensed open
base models and add value through selection, light instruction/format tuning, and convenient
deployment packaging. This overview documents that honestly.

\subsection{What Zen Is---and Is Not}

\begin{itemize}
    \item \textbf{Is}: a curated catalogue of derivatives of open upstream models (mostly Qwen3,
          Apache-2.0), repackaged with deployment artifacts and, in some cases, light fine-tuning.
    \item \textbf{Is not}: a from-scratch pretraining effort, a novel ``Zen MoDE'' architecture, or a
          set of models with independently verified order-of-magnitude efficiency or accuracy
          advantages over their own upstreams.
\end{itemize}

\subsection{Corrections to Prior Versions}

The following previously published claims were inaccurate and are withdrawn:
\begin{itemize}
    \item A proprietary ``Zen MoDE (Mixture of Distilled Experts)'' architecture --- the underlying
          models are standard Qwen3 (dense and Qwen3-30B-A3B MoE) and other open architectures.
    \item Specific homegrown parameter counts including a ``480B'' code model and ``235B'' design
          models --- these do not correspond to the shipped bases.
    \item ``Performance comparable to models 10$\times$ their size'' and ``reducing energy consumption
          by up to 98\%'' --- unsubstantiated marketing claims; removed.
    \item Per-model benchmark tables, environmental-impact (energy/CO$_2$/water) tables, and
          per-device throughput figures --- these were fabricated and are removed in favor of pointers
          to upstream-reported, attributed numbers.
\end{itemize}

\section{Model Family: Bases and Licenses}

\subsection{Provenance Table}

Table~\ref{tab:provenance} lists each Zen language model with its real upstream base and license.
Zen language models are predominantly Qwen3 forks; a small number are based on other open models.
All listed Qwen3 bases are released by Alibaba Cloud under Apache-2.0 \cite{qwen3report, qwen3hf}.

\begin{table}[H]
\centering
\small
\begin{tabular}{lllp{4.2cm}}
\toprule
\textbf{Zen Model} & \textbf{Upstream Base} & \textbf{License} & \textbf{Notes} \\
\midrule
Zen-Nano        & Qwen3-0.6B        & Apache-2.0 & Smallest dense tier; true edge size \\
Zen-Eco         & Qwen3-4B          & Apache-2.0 & Consumer/laptop tier \\
Zen (base)      & Qwen3-8B          & Apache-2.0 & Default instruction model \\
Zen-Pro         & Qwen3-8B          & Apache-2.0 & Reasoning-tuned (long-CoT/RLVR) \\
Zen-Omni        & Qwen3-30B-A3B (MoE) & Apache-2.0 & 30.5B total / 3.3B active MoE \\
Zen-Next        & Qwen3-32B         & Apache-2.0 & Largest dense Qwen3 tier \\
\midrule
Zen (DeepSeek tier) & DeepSeek (e.g.\ V-series) & per upstream & Non-Qwen base; verify upstream license \\
Zen (Granite tier)  & IBM Granite      & Apache-2.0 & Non-Qwen base (IBM Granite) \\
\bottomrule
\end{tabular}
\caption{Zen language models and their real upstream bases. Parameter counts, native context, and
benchmark numbers are those of the named upstream; see the corresponding upstream technical report.
Non-Qwen tiers must carry the license of their specific upstream base.}
\label{tab:provenance}
\end{table}

\paragraph{On vision, image, and speech members.} Zen also ships vision/image and speech models. Those
are likewise derivatives of open upstreams (for example, Qwen vision-language and ASR models, and
permissively licensed image/video generators) and are documented in their own whitepapers with
explicit base-model attribution. A prior license audit found that several Zen media releases had been
mislabeled relative to their true upstreams; those have been corrected or removed in the model
repositories. This overview focuses on the language family; consult each model's individual card and
whitepaper for authoritative per-model provenance and license.

\subsection{The Qwen3 Base Series (for reference)}

For context, the upstream Qwen3 series \cite{qwen3report} provides open-weight dense models at 0.6B,
1.7B, 4B, 8B, 14B, and 32B, plus two mixture-of-experts models, Qwen3-30B-A3B (30.5B total / 3.3B
active) and Qwen3-235B-A22B (235B total / 22B active). All are released under Apache-2.0. The Zen
language tiers select from the smaller end of this series (0.6B/4B/8B/32B dense and the 30B-A3B MoE);
Zen does not ship a from-scratch model larger than its chosen Qwen3 bases.

\section{Architecture (Inherited)}

Zen language models inherit the architecture of their Qwen3 bases without modification. In summary,
the dense Qwen3 models are decoder-only transformers with:
\begin{itemize}
    \item Grouped-query attention (GQA) \cite{ainslie2023gqa};
    \item SwiGLU feed-forward blocks;
    \item RMSNorm normalization;
    \item Rotary position embeddings (RoPE), with long context via YaRN \cite{peng2023yarn};
    \item A 151{,}936-token multilingual BPE vocabulary.
\end{itemize}
Qwen3-30B-A3B additionally uses a sparse mixture-of-experts feed-forward layer (128 experts, 8
activated). There is no proprietary ``MoDE'' routing; the only MoE in the Zen language family is the
upstream Qwen3-30B-A3B design. Per-model hyperparameters (layer counts, hidden sizes, head
configurations, native context) are exactly those published in the Qwen3 technical report and model
cards \cite{qwen3report, qwen3hf}.

\section{What the Zen Team Actually Does}

\subsection{Selection and Curation}

We choose an appropriate open base for each tier (size, license, capability) and document the choice.

\subsection{Light Instruction / Format Tuning}

For chat-tuned variants we may apply light supervised fine-tuning (response-masked loss, low learning
rate, few epochs) and, optionally, preference optimization (DPO) \cite{rafailov2023dpo} to adjust
instruction-following style. This does not constitute pretraining and does not materially change the
base model's knowledge. Where reasoning tuning is applied (e.g., Zen-Pro), it consists of long
chain-of-thought SFT and optional reinforcement learning from verifiable rewards on math/code.

\subsection{Packaging and Deployment Artifacts}

The principal deliverable is convenient, multi-format deployment. We publish, per model:
\begin{itemize}
    \item BF16 SafeTensors reference weights;
    \item GGUF quantizations (e.g., Q4\_K\_M, Q8\_0) for llama.cpp on CPU/GPU;
    \item MLX 4-bit builds for Apple Silicon.
\end{itemize}
Quantization uses standard, widely used post-training schemes applied to the released weights. Any
quality cost of quantization relative to BF16 is measured on the released artifacts and reported
alongside, rather than asserted.

\section{Training Methodology (Upstream)}

Large-scale pretraining of the underlying models was performed by the upstream teams (Alibaba's Qwen
team for Qwen3; the respective teams for any non-Qwen bases). The Zen team does not run pretraining.
For the pretraining corpus, scale, data composition, and pretraining procedure, see the upstream
technical reports---principally the Qwen3 technical report \cite{qwen3report}. Post-training that the
Zen team performs is limited to the light fine-tuning described above and is documented per model.

\section{Evaluation and Benchmarks}

\paragraph{No fabricated tables.} This overview deliberately does not reproduce per-model benchmark
tables. The capabilities of each Zen model are, to first order, those of its upstream base. For
rigorous, independently reproducible benchmark results---MMLU, GSM8K, MATH, HumanEval, multilingual
evaluation, long-context retrieval, and reasoning suites---we refer the reader to the upstream
technical reports and model cards, principally the Qwen3 technical report \cite{qwen3report} and the
per-model Qwen3 cards \cite{qwen3hf}. Where a specific Zen release reports numbers, they are measured
on the released artifact under a stated harness and presented alongside the corresponding upstream
baseline so that the effect of any Zen fine-tuning is transparent.

\paragraph{No fabricated efficiency claims.} The previously stated efficiency advantages
(``10$\times$ smaller,'' ``up to 98\% energy reduction'') and the environmental-impact tables were
not derived from measurements and are removed. The genuine efficiency story is simply that smaller
models in the family (e.g., the 0.6B/4B tiers) are cheaper to run than larger ones, and that 4-bit
quantization reduces memory and increases throughput at a measurable accuracy cost---facts that hold
for the upstream models as well.

\section{Deployment and Integration}

\subsection{Deployment Formats}

\begin{table}[H]
\centering
\small
\begin{tabular}{llp{6.5cm}}
\toprule
\textbf{Format} & \textbf{Typical Precision} & \textbf{Platform / Runtime} \\
\midrule
SafeTensors & BF16/FP16 & Hugging Face Transformers, vLLM (reference weights) \\
GGUF        & Q4\_K\_M, Q8\_0 & llama.cpp (CPU/GPU) \\
MLX         & 4-bit & Apple Silicon \\
ONNX        & INT8 & Cross-platform runtimes \\
\bottomrule
\end{tabular}
\caption{Deployment formats shipped for Zen language models. Memory footprint and throughput depend on
the base model size, quantization, and hardware; we report measured figures per artifact rather than
fixed percentages.}
\end{table}

\subsection{Hardware Footprint}

Memory requirements scale with the base model size and quantization in the usual way (roughly
$\sim$2 bytes/parameter in BF16, $\sim$0.5 bytes/parameter at 4-bit, plus KV cache). For example, an
8B-class model (Zen / Zen-Pro) runs in BF16 on a single 24\,GB GPU and in 4-bit on commodity hardware;
the 0.6B tier (Zen-Nano) is genuinely small enough for edge devices. We do not publish a fixed
memory table here because the values are simply those implied by the chosen Qwen3 base size.

\subsection{Python Integration}
\begin{lstlisting}[language=Python]
# Zen models load via the standard Transformers API,
# because they are Qwen3-architecture models.
from transformers import AutoModelForCausalLM, AutoTokenizer

model = AutoModelForCausalLM.from_pretrained("zenlm/zen-eco-4b-instruct")
tokenizer = AutoTokenizer.from_pretrained("zenlm/zen-eco-4b-instruct")

messages = [{"role": "user", "content": "Solve this problem step by step."}]
inputs = tokenizer.apply_chat_template(messages, return_tensors="pt")
outputs = model.generate(inputs, max_new_tokens=2000)
\end{lstlisting}

\subsection{REST API}
\begin{lstlisting}[language=bash]
# Serve via any Qwen3-compatible OpenAI-style server (e.g., vLLM).
# Query with an OpenAI-compatible request:
curl -X POST http://localhost:8080/v1/completions \
  -H "Content-Type: application/json" \
  -d '{"model": "zen-eco-4b-instruct", "prompt": "Hello, world!", "max_tokens": 100}'
\end{lstlisting}

\section{Licensing and Attribution}

Most Zen language models derive from Qwen3, released by Alibaba Cloud under Apache-2.0, which permits
commercial use, modification, and redistribution subject to attribution and the terms of the license.
Zen distributes these derivatives under the same Apache-2.0 terms with the required upstream
attribution and NOTICE, and records the \texttt{base\_model} for each release. \textbf{Non-Qwen
members must carry the license of their specific upstream base} (for example, IBM Granite is
Apache-2.0; DeepSeek-based members must follow the applicable DeepSeek model license). Relabeling
upstream weights does not extinguish upstream license terms; each model card states its authoritative
license. Users should consult the upstream model cards and LICENSE files for definitive terms.

\section{Safety and Alignment}

Zen models inherit the safety properties and limitations of their upstream bases. Any additional
safety behavior comes from the light instruction/preference tuning described above; we do not claim a
distinct ``Constitutional AI'' training program or specific red-teaming hour counts that were not
performed. Safety alignment is incomplete, as it is for the upstream models, and adversarial prompts
can elicit policy-violating outputs.

\section{Limitations}

Because Zen models are derivatives, they inherit the limitations of their upstreams: factual
hallucination, degraded performance on long multi-step reasoning, incomplete safety alignment, and
the accuracy cost of low-bit quantization. Light fine-tuning can also regress some base-model
capabilities; any claimed improvement should be verified against the upstream baseline. The Zen family
does not provide capabilities beyond what its chosen open bases provide.

\section{Conclusion}

The Zen AI Model Family is best understood as a curated, well-packaged distribution of open upstream
models---predominantly Apache-2.0 Qwen3, with a few non-Qwen members---rather than a from-scratch
model program. Its value lies in selection, light instruction/format tuning, and convenient
multi-format deployment artifacts, all under the upstream licenses with attribution. This corrected
overview removes the prior fabricated architecture, parameter-count, efficiency, and benchmark claims
and points readers to the upstream technical reports for authoritative details. Honest provenance is,
we believe, more useful to adopters than inflated marketing.

\section*{Acknowledgments}

We thank the upstream teams whose open models make the Zen family possible---principally the Qwen team
at Alibaba Cloud for the Apache-2.0 Qwen3 series---and the open-source communities behind Transformers,
llama.cpp/GGML, and MLX.

\appendix

\section{Model Availability}

Zen models and their per-model cards (with \texttt{base\_model} attribution and license) are available
at:
\begin{itemize}
    \item \textbf{HuggingFace}: \url{https://huggingface.co/zenlm}
    \item \textbf{GitHub}: \url{https://github.com/zenlm/zen}
    \item \textbf{Documentation}: \url{https://docs.hanzo.ai/zen}
\end{itemize}

\section{Citation}

\begin{verbatim}
@misc{zenfamily2025,
  title  = {The Zen AI Model Family: An Honest Overview of Finetunes
            and Packaging of Open Models},
  author = {Zen LM Research Team (Hanzo / Lux / Zoo Labs)},
  note   = {Derivatives of open upstream models, principally
            Apache-2.0 Qwen3 (Alibaba Cloud)},
  year   = {2025}
}
\end{verbatim}

\begin{thebibliography}{99}
\bibitem{qwen3report} Qwen Team, Alibaba Cloud, ``Qwen3 Technical Report,'' \textit{arXiv:2505.09388}, 2025.
\bibitem{qwen3hf} Qwen Team, ``Qwen3 model cards and configurations,'' Hugging Face, \url{https://huggingface.co/Qwen}, 2025.
\bibitem{vaswani2017attention} A.~Vaswani et al., ``Attention is All You Need,'' \textit{NeurIPS}, 2017.
\bibitem{ainslie2023gqa} J.~Ainslie et al., ``GQA: Training Generalized Multi-Query Transformer Models,'' \textit{EMNLP}, 2023.
\bibitem{peng2023yarn} B.~Peng et al., ``YaRN: Efficient Context Window Extension of Large Language Models,'' \textit{arXiv:2309.00071}, 2023.
\bibitem{shazeer2017outrageously} N.~Shazeer et al., ``Outrageously Large Neural Networks: The Sparsely-Gated Mixture-of-Experts Layer,'' \textit{ICLR}, 2017.
\bibitem{hu2021lora} E.~J. Hu et al., ``LoRA: Low-Rank Adaptation of Large Language Models,'' \textit{arXiv:2106.09685}, 2021.
\bibitem{dettmers2023qlora} T.~Dettmers et al., ``QLoRA: Efficient Finetuning of Quantized LLMs,'' \textit{NeurIPS}, 2023.
\bibitem{frantar2022gptq} E.~Frantar et al., ``GPTQ: Accurate Post-Training Quantization for Generative Pre-trained Transformers,'' \textit{arXiv:2210.17323}, 2022.
\bibitem{ouyang2022training} L.~Ouyang et al., ``Training language models to follow instructions with human feedback,'' \textit{NeurIPS}, 2022.
\bibitem{rafailov2023direct} R.~Rafailov et al., ``Direct Preference Optimization,'' \textit{NeurIPS}, 2023.
\bibitem{rafailov2023dpo} R.~Rafailov et al., ``Direct Preference Optimization: Your Language Model is Secretly a Reward Model,'' \textit{NeurIPS}, 2023.
\end{thebibliography}

\end{document}
